\documentclass[12pt]{article}
\usepackage[margin=1in]{geometry}
\usepackage{enumitem}
\usepackage{titlesec}

\title{Project Overview: Automated Index Generation}
\author{}
\date{}

\begin{document}

\maketitle

The goal was to modify the PostgreSQL executor to track sequential scans on table columns and automatically trigger a CREATE INDEX command once a specific frequency threshold is reached.

\section*{Changes Made to the Codebase}

\subsection*{1. src/include/utils/autoindex.h}

\textbf{Changes:}
\begin{itemize}[leftmargin=*]
\item Defined the \texttt{AutoIndexKey} struct (mapping relid and attnum).
\item Defined the \texttt{AutoIndexEntry} struct (storing the scan\_count).
\item Set the \texttt{AUTOINDEX\_THRESHOLD} (the trigger limit, currently set to 5).
\item Declared function prototypes for shared memory initialization and tracking.
\end{itemize}

\noindent\textbf{Reasoning:}

\noindent Centralizing these definitions ensures that both the shared memory management and the executor-level tracking use the same data structures and threshold logic.

\section*{2. src/backend/storage/ipc/ipci.c}

\textbf{Changes:}
\begin{itemize}[leftmargin=*]
\item Modified \texttt{CreateSharedMemoryAndSemaphores} to include calls to \texttt{AutoIndexShmemSize() and AutoIndexShmemInit()}.
\end{itemize}

\noindent\textbf{Reasoning:}

\noindent PostgreSQL uses a fixed-size shared memory segment. We must register our new hash table during the server's pre-initialization phase so that all background processes (backends) share the same scan counters.

\section*{3. src/backend/executor/nodeSeqscan.c}

\textbf{Changes:}
\begin{itemize}[leftmargin=*]
\item Added a hook inside ExecInitSeqScan.
\item The code iterates through the scan's filter predicates (plan->qual) to identify equality expressions (e.g., \texttt{WHERE col = 100}).
\item Extracts the AttrNumber and calls TrackEqualityPredicate().
\end{itemize}

\noindent\textbf{Reasoning:}

\noindent ExecInitSeqScan is the "entry point" for every sequential scan. By intercepting the plan here, we can accurately record when a developer is querying an unindexed column repeatedly.

\section*{4. src/backend/utils/misc/autoindex.c}

This is the core logic file. Significant refinements were made here to handle PostgreSQL’s strict memory and transaction safety rules.

\subsection*{Key Logic: TrackEqualityPredicate}

\begin{itemize}[leftmargin=*]
\item \textbf{System Filter:} Added an \texttt{OID} check (\texttt{relid < 16384}) to ignore System Catalogs.

```
Why: This prevents the recursive loop where the auto-indexer tries to index Postgres's own metadata tables during a \texttt{\textbackslash d} command.

\item \textbf{Shared Memory Access:} Uses LWLockAcquire (AddinShmemInitLock) to ensure thread-safe increments of scan counts across multiple concurrent users.
```

\end{itemize}

\subsection*{Key Logic: ExecuteAutoIndex}

\begin{itemize}[leftmargin=*]
\item \textbf{Variable Scoping:} All declarations were moved to the top of the function to comply with ISO C89 standards.

```
\item \textbf{Memory Management:} Added pstrdup() for relation and attribute names.

Why: This copies strings into local memory, preventing "use-after-free" errors if the system cache is invalidated during index creation.

\item \textbf{Manual Node Construction:} Instead of using the version-dependent makeIndexElem(), we manually initialized the IndexElem struct.

Why: This bypassed "undeclared function" errors caused by refactored headers in the PostgreSQL source.

\item \textbf{Transaction Safety (\texttt{PG\_TRY / PG\_CATCH}):} Wrapped the DefineIndex call in a trial block.

Why: This allows the original \texttt{SELECT} query to finish even if the index creation hits a lock conflict, preventing a full session crash.

\item \textbf{Concurrency Setting:} Set \texttt{stmt->concurrent = false}.

Why: While concurrent builds are great for production, they cannot be triggered from inside an active transaction. Standard builds are significantly more stable for this specific executor-hook implementation.

\end{itemize}

\section*{Final Results}

The system now successfully:
\begin{itemize}[leftmargin=*]
\item Detects repeated scans on user tables.
\item Ignores internal system queries.
\item Automatically generates valid B-Tree indices.
\item Maintains database stability by catching resource-owner conflicts.
\end{itemize}

\section*{Current Implementation Status}

\begin{itemize}[leftmargin=*]
\item \textbf{Success:} Indices are appearing in \texttt{\textbackslash d} after reaching the threshold.
\end{itemize}

\section*{Future Extensions}

\textbf{1. Cost-Benefit Thresholding}

\begin{itemize}
    \item \textbf{Concept:} Replace the "magic number" threshold with a financial model of system resources.

    \item \textbf{Logic:} Trigger creation only when
    \texttt{(N $\times$ Savings) > Creation Cost} where \texttt{Savings = cost\_seqscan - cost\_indexscan}

    \item \textbf{Implementation:} Utilize internal cost\_seqscan and cost\_indexscan metrics to estimate the "break-even" point for CPU and I/O.
\end{itemize}

\noindent\textbf{2. Asynchronous Background Worker (BGW)}

\begin{itemize}
    \item \textbf{Concept:} Decouple index creation from the user's \texttt{SELECT} process.

    \item \textbf{Logic:} The Executor places a request in a shared memory queue; a dedicated \texttt{BGW} process handles the \texttt{DDL}.

    \item \textbf{Benefit:} Eliminates "Resource Owner" conflicts and session latency. Enables the safe use of \texttt{CREATE INDEX CONCURRENTLY}.
\end{itemize}

\noindent\textbf{3. Selectivity Analysis}

\begin{itemize}
    \item \textbf{Concept:} Prevent "weak" indexes on low-cardinality columns (e.g., Booleans, Gender).

    \item \textbf{Logic:} Query pg\_statistic to check the stadistinct value.

    \item \textbf{Action:} Abort auto-indexing if the column filters out less than a significant percentage (e.g., \texttt{<20\%}) of the total table rows.
\end{itemize}

\noindent\textbf{4. Composite (Multi-Column) Indexing}

\begin{itemize}
    \item \textbf{Concept:} Recognize patterns of multi-clause queries.

    \item \textbf{Logic:} If \texttt{WHERE A = x AND B = y} is common, detect the co-occurrence of predicates in the plan tree.

    \item \textbf{Outcome:} Generate a single composite index \texttt{(A, B)} instead of two redundant single-column indexes.
\end{itemize}

\noindent\textbf{5. "Human-in-the-Loop" Approval Mode}

\begin{itemize}
    \item \textbf{Concept:} Provide administrative oversight for database changes.

    \item \textbf{Logic:} Instead of immediate execution, log proposed indexes to a system table (e.g., pg\_index\_proposals).

    \item \textbf{Outcome:} Allows DBAs to review, test, and "approve" auto-generated suggestions before they impact production storage.
\end{itemize}

\end{document}
